\chapter*{Executive Summary}
\addcontentsline{toc}{chapter}{Executive Summary}

\textbf{Problem.} Neighbours lend to one another, but a loan ends with an item handed to someone the lender may not know. Existing platforms either show one averaged rating, which hides how little evidence stands behind it, or show nothing at all. Neighbourhood evidence is sparse (most members have few loans), role-dependent (a good lender is not necessarily a good borrower), indirect (a borrower may be known to the lender's neighbours but not to the lender) and manipulable (a small group can rate itself up). Digital Mohalla treats the question \textit{``is there enough trustworthy evidence for this loan?''} as its complex computing problem.

\textbf{Response to the proposal defence.} The panel approved the proposal with minor modifications and made three points: show that location is a real issue, the platform has positive social potential, and the technical challenge looked limited. The third point changed the project most. The team narrowed the first iteration to the assurance engine and the transaction lifecycle that feeds it, and answered the location point with a comparison of five locality methods, a threat model and eighteen test checks on the chosen rule (Chapter~\ref{ch:proposal-feedback-iteration}).

\textbf{What Iteration 1 delivered.} All three planned slices are complete and merged:
\begin{itemize}
\item \textbf{Slice A, workflow skeleton:} an append-only, hash-chained event store; a 14-state borrowing lifecycle whose transitions are checked in Python and again by a PostgreSQL foreign key; double booking prevented by a database exclusion constraint; login and a REST API; a resident app; CI and CD.
\item \textbf{Slice B, assurance core:} outcome extraction from closed loans, a recency-weighted Beta posterior with a 90\% credible interval and similarity weighting, trust edges and bounded $k$-hop search ($k=3$), a decision policy of rules R1--R9 (R6 in two variants) with safeguards and plain-language explanations, a scenario suite in CI and the lender's evidence screen.
\item \textbf{Slice C, human loop:} partner-clustering and review-similarity signals with a retaliation check, disputes that freeze a loan, moderator decisions on disputes and holds, privacy checks on every API response and log line, request fuzzing, and the moderator app.
\end{itemize}

\textbf{Evidence.} 623 backend tests (unit, integration on PostgreSQL 16, property-based and scenario tests) and 29 frontend tests pass in CI, with 97\% line coverage of the backend code, and every pull request runs them before merge. All seven scenarios required at mid-term (ES-01 to ES-06 and ES-13) return the documented state \textit{and} explanation, against a target of seven of thirteen (NFR-Q-01); eleven of the thirteen reproduce their expected state end to end on synthetic data. At $k=3$ the indirect-evidence search takes 4.5~ms at the 95th percentile on a 2,026-user synthetic graph and finds a path for 93.2\% of loan pairs. The API and the web app are deployed on Render with a Neon PostgreSQL database, from the \code{main} branch after CI passes.

\textbf{Honest limits.} All parameters (half-life, thresholds, decay, signal weights) are initial hypotheses checked only on synthetic data. The locality method is designed but not built; photo uploads, the audit-log screen, three manipulation signals (S2, S3, S5) and the scheduler move to Iteration 2. Generative AI (Claude Code) drafted most first-draft code under the team's direction; its use is declared per module in Section~\ref{sec:tools-tool-assisted-development}.

\textbf{Next.} Iteration 2 completes the core modules (membership and locality, the remaining lifecycle timers, the remaining signals and reviews, case information requests and appeals) and builds the simulator, baselines and FYP-2 evaluation plan that will test whether the engine's uncertainty actually improves decisions over simple averages.
